\section{Discussion}
\label{sec:discussion}

\subsection{Empirical Associations with Radar-Topographic Geometry}
\label{sec:disc:physics}

The empirical performance of TopoRadar-Net relative to standard attention-augmented and Siamese convolutional baselines is consistent with known coupling between microwave backscatter and steep mountain topography. Conventional computer vision models treat SAR intensity channels as flat 2D arrays, whereas TopoRadar-Net additionally receives terrain and viewing-geometry coordinates. The reported ablations and strata quantify sensitivity to those inputs; they do not isolate a causal internal mechanism.

As demonstrated by Eq.~(\ref{eq:lia_physics}), local incidence angle $\psi$ is governed by terrain slope $\theta$, aspect azimuth $\alpha$, and satellite radar look azimuth $\phi_{\text{look}}$. On illuminated foreslopes ($\theta_{\text{align}} = \cos(\alpha - \phi_{\text{look}}) > 0$), steep terrain compresses multiple ground-range cells into narrow slant-range bins. This foreshortening can elevate radar backscatter and generate bright linear ridges. RTCAB makes local-incidence and aspect-alignment features available to the SAR representation; the observed region- and stratum-specific changes are compatible with, but do not directly demonstrate, learned foreslope suppression.

Conversely, on backslopes facing away from the radar ($\theta_{\text{align}} < 0, \psi > 50^\circ$), undisturbed smooth snow surfaces can appear radiometrically dark due to specular reflection away from the radar antenna. Avalanche debris may then exhibit strong diffuse surface and volume backscatter ($\Delta \sigma^0 > 0$). TopoRadar-Net records higher Seed-42 F1 than Attention U-Net in the reported backslope strata, but activation-level evidence would be required to attribute that difference specifically to dynamic spatial gating.

Geo-Loss encodes a heuristic slope prior rather than a physical-validity constraint. Its mask overlaps genuine positive labels, most strongly in the Nuuk 2016 training event (28.79\% of valid positives), and the held-out subgroup audit records lower positive recall than No-GeoLoss in both test systems. The aggregate Troms{\o} precision/F1 difference and the subgroup recall penalty therefore describe a trade-off; without a matched generic regularizer they do not establish a terrain-specific mechanism.

\subsection{Cross-System Transferability Across Global Mountain Regimes}
\label{sec:disc:transfer}

A primary challenge in machine learning for natural hazards is model overfitting to localized geographic and environmental training conditions. The AvalCD benchmark spans four distinct mountain systems characterized by starkly contrasting climatic, snowpack, and geomorphic regimes:
\begin{enumerate}
    \item \textbf{Temperate Alpine Regime (European Alps):} High recreation density, moderate elevation relief, and mixed maritime-continental snowpacks characterized by distinct layer metamorphism and frequent spring wet-snow transitions.
    \item \textbf{High-Altitude Continental Regime (Pamir Mountains, Central Asia):} Extreme vertical relief ($>2,000$~m descent paths), hyper-arid winter climate, deeply facetted depth hoar, and massive powder/slab avalanche cycles capable of burying entire valleys.
    \item \textbf{Maritime Sub-Arctic Regime (Troms{\o}, Norway):} Coastal mountains with heavy, wind-scoured maritime snowpacks, frequent rain-on-snow events, and polar darkness throughout mid-winter.
    \item \textbf{Polar Fjord Alpine Regime (Nuuk, Greenland):} Sub-polar maritime conditions with persistent high-velocity winds and coastal fjord relief.
\end{enumerate}

The model was trained on Alps and Greenland events, selected on Alps and Greenland validation events, and evaluated without site-specific tuning on held-out Troms{\o} and Pamir scenes. It records $77.40\% \pm 1.86\%$ and $50.18\% \pm 4.40\%$ F1, respectively. These observed held-out results change climate, relief, acquisition geometry, and avalanche population simultaneously and therefore do not identify physical coordinates as the cause of transfer.

\subsection{Operational Implications for Alpine Civil Protection and Search-and-Rescue}
\label{sec:disc:operational}

For alpine emergency managers, highway maintenance authorities, and search-and-rescue teams, satellite-based avalanche detection is only valuable if it provides reliable, actionable intelligence during crisis cycles \citep{eckerstorfer2019near, eaws2022size, techel2018spatial}. Optical remote sensing is consistently thwarted during major blizzard events when rapid avalanche detection is most urgently needed. Sentinel-1 C-band SAR provides the only freely accessible, all-weather observation platform capable of continuous winter surveillance.

Large avalanches can block transportation arteries and bury infrastructure, motivating size-stratified evaluation. On the single held-out Troms{\o} inventory, TopoRadar-Net records mean hit rates of $84.0\%$ for D3 and $95.8\%$ for D4. These retrospective scene-bounded rates neither guarantee identification of future events nor measure civil-protection outcomes.

\begin{table*}[t]
\caption{Conceptual decision-support framework and false-alarm footprint on the evaluated finite-raster valid masks (Troms{\o}: $231.151~\text{km}^2$; Pamir: $350.338~\text{km}^2$ on the nominal 10 m grid). The framework is illustrative and unvalidated for automated dispatch; affine areas scale proportionally with identical ha/$\text{km}^2$ rates.}
\label{tab:operational_triage}
\centering
\footnotesize
\setlength{\tabcolsep}{5pt}
\begin{tabular*}{\textwidth}{@{\extracolsep{\fill}}lp{2.2cm}p{2.3cm}p{4.7cm}p{3.4cm}@{}}
\toprule
\textbf{Triage Level} & \textbf{Score ($p$)} & \textbf{Debris Footprint} & \textbf{Illustrative Hazard Tier} & \textbf{Target Stakeholder} \\
\midrule
\textbf{Tier 1} & $p < 0.40$ & $< 0.1 \text{ ha}$ ($< 10 \text{ px}$) & Low confidence; routine background monitoring & Warning Services (SLF, NVE) \\
\textbf{Tier 2} & $0.40 \le p < 0.70$ & $0.1\text{--}1.0 \text{ ha}$ ($10\text{--}100 \text{ px}$) & Moderate confidence; recommended expert review or optical cueing & Regional Forecasters \\
\textbf{Tier 3} & $p \ge 0.70$ & $\ge 1.0 \text{ ha}$ (Major debris) & High confidence; prioritized inspection along transport corridors & Civil Protection / Road Auth. \\
\midrule
\multicolumn{5}{l}{\textbf{Regional False-Alarm Footprint (Valid Area: Troms{\o} $231.151 \text{ km}^2$, Pamir $350.338 \text{ km}^2$)}} \\
\midrule
\textbf{Region} & \textbf{Method} & \textbf{Detected Debris (ha)} & \textbf{False-Alarm Area (ha / $\text{km}^2$)} & \textbf{False-Alarm Rate (ha / $\text{km}^2$)} \\
\midrule
\multirow{3}{*}{Scandinavian Arctic (Troms{\o})} & SiamUNet-diff & $218.33 \text{ ha}$ & $43.13 \text{ ha}$ ($0.43 \text{ km}^2$) & $0.19 \text{ ha} / \text{km}^2$ (Misses $121.36 \text{ ha}$) \\
 & Attention U-Net & $261.49 \text{ ha}$ & $68.34 \text{ ha}$ ($0.68 \text{ km}^2$) & $0.30 \text{ ha} / \text{km}^2$ \\
 & \textbf{TopoRadar-Net (Ours)} & $\mathbf{270.09 \text{ ha}}$ & $\mathbf{87.82 \text{ ha}}$ ($\mathbf{0.88 \text{ km}^2}$) & $\mathbf{0.38 \text{ ha} / \text{km}^2}$ (Detects $95.8\%$ D4) \\
\midrule
\multirow{3}{*}{Pamir Mountains (Pish)} & SiamUNet-conc & $163.76 \text{ ha}$ & $302.63 \text{ ha}$ ($3.03 \text{ km}^2$) & $0.86 \text{ ha} / \text{km}^2$ (Severe Clutter) \\
 & Attention U-Net & $150.43 \text{ ha}$ & $242.77 \text{ ha}$ ($2.43 \text{ km}^2$) & $0.69 \text{ ha} / \text{km}^2$ \\
 & \textbf{TopoRadar-Net (Ours)} & $\mathbf{130.25 \text{ ha}}$ & $\mathbf{169.89 \text{ ha}}$ ($\mathbf{1.70 \text{ km}^2}$) & $\mathbf{0.48 \text{ ha} / \text{km}^2}$ ($\mathbf{30.0\%}$ Clutter Reduction) \\
\bottomrule
\end{tabular*}
\end{table*}

Table~\ref{tab:operational_triage} summarizes the quantitative clutter footprint across mountain regimes and outlines an illustrative conceptual triage framework:
\begin{enumerate}
    \item \textbf{Physical-Unit Footprints:} On the Pamir valid mask ($350.338~\text{km}^2$ nominal), TopoRadar-Net records $169.89$ ha of false alarms ($0.48~\text{ha/km}^2$), $72.88$ ha ($30.0\%$) below Attention U-Net and $132.73$ ha ($43.9\%$) below SiamUNet-conc; affine physical differences are $55.63$ and $101.33$ ha. In Troms{\o}, its density is $0.38~\text{ha/km}^2$ and its false-alarm area exceeds Attention U-Net, so the trade-off is region dependent.
    \item \textbf{Calibrated Corridor Proxy:} We fit Platt calibration and an F2-weighted component threshold only on the two validation scenes, then intersect $\ge0.1$ ha test components with a frozen OpenStreetMap major-road snapshot (30 m buffer) \citep{openstreetmap2026}. In Pamir (8 road-intersecting ground-truth components), TopoRadar-Net achieves $83.0\% \pm 2.3\%$ alert precision and $62.5\% \pm 10.2\%$ recall; Attention U-Net achieves $77.4\% \pm 6.1\%$ and $70.8\% \pm 5.9\%$. Both models miss the single Troms{\o} corridor component in all seeds. TopoRadar Brier score improves from $0.282$ to $0.208$ in Troms{\o} and $0.462$ to $0.198$ in Pamir after calibration. This post-hoc proxy falsifies automated-dispatch readiness while quantifying a bounded Pamir use case.
    \item \textbf{Cached End-to-End Profile:} With the OSM snapshot cached, all measured TopoRadar-Net runs for raster loading, road-buffer rasterization, full-scene inference, calibration, connected components, and corridor intersection remain below 5 s per scene on Apple MPS. Wall-clock values vary with device warm state and are verified by bounds rather than byte identity. The profile excludes data acquisition and human review and therefore does not establish response latency or decision benefit.
\end{enumerate}

\subsection{Limitations and Future Research}
\label{sec:disc:limitations}

Several physical and technical limitations warrant disclosure:
\begin{itemize}
    \item \textbf{Augmentation Convention:} The released checkpoints use the disclosed checkpoint-era convention in which aspect sine/cosine components are sign-adjusted under reflections while the precomputed orbital look-alignment channel is spatially transformed without directional recomputation; orthogonal rotations similarly preserve channel values. The resulting training tensor is not fully equivariant under these transformations. Consequently, the directional results are interpreted as empirical sensitivity rather than a causal validation of the manifold representation. A future confirmatory rerun should compare a fully equivariant augmentation against the released convention.
    \item \textbf{Heuristic Slope Prior:} Geo-Loss penalizes positive predictions whenever slope is below $5^\circ$ or above $65^\circ$, including genuine mapped debris. The conflict is concentrated in one training event but is present in every split; held-out subgroup recall is lower than No-GeoLoss. Future work should compare label-aware and generic regularizers on a fresh protected cohort.
    \item \textbf{DEM Quality and Geometric Resolution:} TopoRadar-Net relies on the Copernicus GLO-30 DEM. While GLO-30 provides global coverage, localized vertical errors or smoothing over sharp knife-edge ar{\^e}tes can propagate inaccuracies into the local incidence angle calculation. In regions with available high-resolution airborne LiDAR or stereo-photogrammetric DEMs (1--5~m), topographic conditioning can be further refined.
    \item \textbf{Dielectric Attenuation in Wet Snow Cycles:} Liquid water changes microwave permittivity and can reduce C-band contrast during rain-on-snow events or warm spells. Bitemporal differencing may therefore mask debris or create broad absorption anomalies. Future work should integrate multi-temporal radiometry or polarimetric decomposition parameters.
    \item \textbf{Temporal Revisit Constraints:} A several-day repeat interval can aggregate multiple releases into one composite deposit. Sentinel-1C and Sentinel-1D launched on 5 December 2024 and 4 November 2025, respectively \citep{esa2026sentinel1}; the expanded constellation improves continuity but does not remove acquisition-latency or event-timing ambiguity.
    \item \textbf{Operational Proxy, Not Stakeholder Validation:} The calibrated corridor analysis uses a committed OpenStreetMap snapshot and retrospective labels. It measures component calibration, route proximity, and cached processing time, not dispatch quality, road closure decisions, or stakeholder benefit. The Troms{\o} miss and the small number of corridor-intersecting avalanches preclude deployment claims.
\end{itemize}
